\documentclass[10pt,a4paper,titlepage]{jreport} % ドキュメントクラスを1つに統一


\usepackage[top=25truemm,bottom=25truemm,left=20truemm,right=20truemm]{geometry} % 余白設定
\usepackage[dvipdfmx]{graphicx} % 画像の挿入
\usepackage{graphicx}
\usepackage{listings} % コードリスト用
\usepackage{jlisting} % 日本語対応のリスト用
\usepackage{jverb} % 日本語のベタ打ち
\usepackage{booktabs}
\usepackage{pgfplots} % グラフ描画用パッケージ
\pgfplotsset{compat=1.18} % バージョン指定
\usepackage{float}
\parindent = 0pt % 段落の字下げを無効化

\usepackage{titlesec}
\usepackage{etoolbox}
\makeatletter
\patchcmd{\chapter}{\if@openleft\cleardoublepage\else\if@openright\cleardoublepage\else\clearpage\fi\fi}{}{}{}
\makeatother
\lstset{breaklines=true, postbreak=\mbox{$\hookrightarrow$}\space, keepspaces=true, escapeinside={\%*}{*)}}

\titleformat{\chapter}[hang] % 章のフォーマットを変更
  {\normalfont\huge\bfseries} % フォント設定
  {\thechapter} % 番号部分の表示形式
  {1em} % 番号とタイトルの間のスペース
  {} % タイトルの前に入れる内容（ここでは空）
\usepackage{xcolor}
\usepackage{amsmath}
\usepackage{amssymb} % ←これで \mathbb{R} が使える

\lstset{
  language=Matlab,              % 言語はMATLAB
  basicstyle=\ttfamily\small,   % フォントは等幅、サイズ小さめ
  keywordstyle=\color{blue},    % キーワード（function, endとか）を青
  commentstyle=\color{green!50!black}, % コメントを深緑
  stringstyle=\color{red!70!black},    % 文字列（'...'）を赤系
  numbers=left,                 % 行番号を左に表示
  numberstyle=\tiny\color{gray},% 行番号をグレーで小さく
  stepnumber=1,                 % 毎行行番号つける
  numbersep=10pt,               % コードとの間隔
  backgroundcolor=\color{gray!10}, % 背景うっすらグレー
  frame=single,                 % 枠をつける
  rulecolor=\color{black},      % 枠線は黒
  breaklines=true,              % 長い行を折り返す
  breakatwhitespace=false,      % スペースでだけ折り返すか（falseにして自然な折り返し）
  showspaces=false,             % スペースは特別に表示しない
  showstringspaces=false,       % 文字列中のスペースも普通に表示
  tabsize=2                     % タブ幅2
}


\title{数値計算：課題10} % タイトル
\author{
  学生番号：242C2016  氏名：奥村直 \\
  \\
  知的システム工学科システム制御コース
  } % 著者
\date{\today} % 日付

\begin{document}

\maketitle

\section{問題の前提と条件}

次の初期値問題を考える。
\begin{equation}
y'(x) = -y(x), \quad 0 \le x \le 1, \qquad y(0)=1
\end{equation}

この微分方程式の解は
\begin{equation}
y(x) = e^{-x}
\end{equation}
であり，
\begin{equation}
y(1) = e^{-1} \approx 0.367879
\end{equation}
である。

% ==================================================
\section{問１：オイラー法}

\subsection{漸化式の導出}

オイラー法は
\begin{equation}
y_{n+1} = y_n + h f(x_n, y_n)
\end{equation}
で定義される。

この問いでは
\[
f(x,y)=-y
\]
であるから，
\begin{equation}
y_{n+1} = y_n - h y_n = (1-h)y_n
\end{equation}
となる。

初期条件は
\[
x_n = nh, \quad y_0 = 1
\]
である。

% ==================================================
\subsection{問２：ソースコード（C言語）オイラー法}

以下にオイラー法を用いたプログラムを示す。

\begin{lstlisting}
#include <stdio.h>
#include <math.h>

int main(void)
{
    double h = 0.1;   // 刻み幅（0.01に変更可）
    double x = 0.0;
    double y = 1.0;
    int N = (int)(1.0 / h);

    for (int n = 1; n <= N; n++) {
        y = y - h * y;
        x = n * h;
        printf("%f %f\n", x, y);
    }

    double exact = exp(-1.0);
    double rel_error = fabs(y - exact) / exact * 100.0;

    printf("y(1) approx = %f\n", y);
    printf("relative error = %f %%\n", rel_error);

    return 0;
}
\end{lstlisting}

% ==================================================
\subsection{数値結果}

刻み幅を
\[
h=0.1,\;0.01
\]
とした場合の $x=1$ における近似解と相対誤差を求めた。

\begin{itemize}
\item $h=0.1$：
\[
y(1)\approx0.348678,\quad \text{相対誤差}\approx5.22\%
\]
\item $h=0.01$：
\[
y(1)\approx0.366032,\quad \text{相対誤差}\approx0.50\%
\]
\end{itemize}

刻み幅を小さくすると誤差が減少することが分かる。

% ==================================================
\section{問３：ホイン法}

\subsection{計算式}

ホイン法は
\begin{align}
k_1 &= f(x_n,y_n) \\
k_2 &= f(x_n+h, y_n + hk_1) \\
y_{n+1} &= y_n + \frac{h}{2}(k_1+k_2)
\end{align}
で定義される。

本問題では
\[
f(x,y)=-y
\]
である。

% ==================================================
\subsection{ソースコード（C言語）：ホイン法}

\begin{lstlisting}
#include <stdio.h>
#include <math.h>

int main(void)
{
    double h = 0.1;   // 刻み幅（0.01に変更可）
    double x = 0.0;
    double y = 1.0;
    int N = (int)(1.0 / h);

    for (int n = 1; n <= N; n++) {
        double k1 = -y;
        double k2 = -(y + h * k1);
        y = y + h * (k1 + k2) / 2.0;
        x = n * h;
        printf("%f %f\n", x, y);
    }

    double exact = exp(-1.0);
    double rel_error = fabs(y - exact) / exact * 100.0;

    printf("y(1) approx = %f\n", y);
    printf("relative error = %f %%\n", rel_error);

    return 0;
}
\end{lstlisting}

% ==================================================
\subsection{考察}

ソースコードを実行した結果，同一の刻み幅においてホイン法からはオイラー法よりも高精度な近似解を得ることができた。
これはホイン法が2次精度を持つ数値解法であるのに対し、
オイラー法が1次精度であることによるものと考えられる。

% ==================================================
\section{総括}

この課題では，オイラー法およびホイン法を用いて常微分方程式の数値解を求めた。
刻み幅を小さくすることで精度が向上し，
またホイン法の方がオイラー法よりも優れた精度を示すことが確認できた。

\end{document}
